\documentclass[11pt,a4paper]{article}

\usepackage[utf8]{inputenc}
\usepackage[T1]{fontenc}
\usepackage[french]{babel}
\usepackage[a4paper,margin=2.2cm]{geometry}
\usepackage{amsmath,amssymb,amsthm}
\usepackage{lmodern}
\usepackage{enumitem}
\usepackage{hyperref}
\usepackage{fancyhdr}
\usepackage{mathrsfs}

\hypersetup{
    colorlinks=true,
    linkcolor=black,
    urlcolor=blue
}

\setlength{\parindent}{0pt}
\setlength{\parskip}{0.6em}

\pagestyle{fancy}
\fancyhf{}
\fancyhead[L]{Programme de colle --- Nombres complexes}
\fancyhead[R]{Terminale générale}
\fancyfoot[C]{\thepage}

\newtheorem{prop}{Démonstration}
\newtheorem*{idee}{Idée}

\newcommand{\C}{\mathbb{C}}
\newcommand{\R}{\mathbb{R}}
\newcommand{\ii}{\mathrm{i}}
\newcommand{\e}{\mathrm{e}}
\newcommand{\Reel}{\operatorname{Re}}
\newcommand{\Imag}{\operatorname{Im}}

\begin{document}

\begin{center}
    {\LARGE \textbf{Programme de colle --- Nombres complexes}}\\[0.5em]
    {\large 15 démonstrations à apprendre}
\end{center}

\vspace{1em}

\hrule
\vspace{1em}

\textbf{Convention.} On note un complexe sous la forme $z=a+\ii b$ avec $a,b\in\R$, son conjugué $\overline z$, et son module $|z|$.

\vspace{1em}
\hrule
\vspace{1em}

\section*{1. Le conjugué d'un produit est le produit des conjugués}

\textbf{Énoncé.} Pour tous $z,z'\in\C$,
\[
\overline{zz'}=\overline z\,\overline{z'}.
\]

\textbf{Type de raisonnement.} Calcul direct par écriture algébrique.

\begin{prop}
Écrivons
\[
z=a+\ii b,\qquad z'=c+\ii d,
\]
avec $a,b,c,d\in\R$.

On calcule d'abord le produit :
\[
zz'=(a+\ii b)(c+\ii d)=ac-bd+\ii(ad+bc).
\]
Par définition du conjugué,
\[
\overline{zz'}=ac-bd-\ii(ad+bc).
\]

Calculons maintenant le produit des conjugués :
\[
\overline z\,\overline{z'}=(a-\ii b)(c-\ii d).
\]
En développant,
\[
(a-\ii b)(c-\ii d)=ac-a\ii d-\ii bc+(-\ii)(-\ii)bd.
\]
Or
\[
(-\ii)(-\ii)=\ii^2=-1.
\]
Ainsi
\[
\overline z\,\overline{z'}=ac-bd-\ii(ad+bc).
\]

Les deux expressions obtenues étant identiques, on conclut que
\[
\overline{zz'}=\overline z\,\overline{z'}.
\]
\end{prop}

\section*{2. Expression du module : \texorpdfstring{$z\overline z=|z|^2$}{zconj(z)=|z|²}}

\textbf{Énoncé.} Pour tout $z\in\C$,
\[
z\overline z=|z|^2.
\]

\textbf{Type de raisonnement.} Calcul direct à partir de l'écriture algébrique et de la définition du module.

\begin{prop}
Soit
\[
z=a+\ii b.
\]
Alors
\[
\overline z=a-\ii b.
\]
On calcule :
\[
z\overline z=(a+\ii b)(a-\ii b)=a^2-(\ii b)^2.
\]
Or
\[
(\ii b)^2=\ii^2b^2=-b^2.
\]
Donc
\[
z\overline z=a^2+b^2.
\]

Par définition du module,
\[
|z|=\sqrt{a^2+b^2}.
\]
En élevant au carré, on obtient
\[
|z|^2=a^2+b^2.
\]
Par conséquent,
\[
z\overline z=|z|^2.
\]
\end{prop}

\section*{3. Formules donnant la partie réelle et la partie imaginaire}

\textbf{Énoncé.} Pour tout $z\in\C$,
\[
\Reel(z)=\frac{z+\overline z}{2},
\qquad
\Imag(z)=\frac{z-\overline z}{2\ii}.
\]

\textbf{Type de raisonnement.} Calcul algébrique par somme et différence de $z$ et $\overline z$.

\begin{prop}
Posons
\[
z=a+\ii b.
\]
Alors
\[
\overline z=a-\ii b.
\]

En additionnant $z$ et $\overline z$, on obtient
\[
z+\overline z=(a+\ii b)+(a-\ii b)=2a.
\]
Ainsi
\[
\frac{z+\overline z}{2}=a=\Reel(z).
\]

En soustrayant $\overline z$ à $z$, on obtient
\[
z-\overline z=(a+\ii b)-(a-\ii b)=2\ii b.
\]
Donc
\[
\frac{z-\overline z}{2\ii}=b=\Imag(z).
\]

On a donc bien
\[
\Reel(z)=\frac{z+\overline z}{2},
\qquad
\Imag(z)=\frac{z-\overline z}{2\ii}.
\]
\end{prop}

\section*{4. Le module d'un produit est le produit des modules}

\textbf{Énoncé.} Pour tous $z,z'\in\C$,
\[
|zz'|=|z|\,|z'|.
\]

\textbf{Type de raisonnement.} Déduction algébrique à partir de l'identité $z\overline z=|z|^2$.

\begin{prop}
On calcule le carré du module de $zz'$ :
\[
|zz'|^2=(zz')\overline{(zz')}.
\]
D'après la propriété du conjugué d'un produit,
\[
\overline{(zz')}=\overline z\,\overline{z'}.
\]
Ainsi,
\[
|zz'|^2=(zz')(\overline z\,\overline{z'})=(z\overline z)(z'\overline{z'}).
\]
Or
\[
z\overline z=|z|^2
\qquad \text{et} \qquad
z'\overline{z'}=|z'|^2.
\]
Donc
\[
|zz'|^2=|z|^2|z'|^2.
\]

Les quantités $|zz'|$ et $|z|\,|z'|$ étant positives ou nulles, l'égalité de leurs carrés entraîne l'égalité :
\[
|zz'|=|z|\,|z'|.
\]
\end{prop}

\section*{5. Formule du module d'un quotient}

\textbf{Énoncé.} Pour tous $z,z'\in\C$ avec $z'\neq 0$,
\[
\left|\frac{z}{z'}\right|=\frac{|z|}{|z'|}.
\]

\textbf{Type de raisonnement.} Déduction à partir de la multiplicativité du module.

\begin{prop}
Comme $z'\neq 0$, le quotient $\dfrac{z}{z'}$ est bien défini et l'on peut écrire
\[
z=\frac{z}{z'}\cdot z'.
\]
En prenant les modules, on obtient
\[
|z|=\left|\frac{z}{z'}\cdot z'\right|.
\]
D'après la propriété de multiplicativité du module,
\[
|z|=\left|\frac{z}{z'}\right|\,|z'|.
\]
Comme $z'\neq 0$, on a $|z'|>0$. On peut donc diviser par $|z'|$, ce qui donne
\[
\left|\frac{z}{z'}\right|=\frac{|z|}{|z'|}.
\]
\end{prop}

\section*{6. Inverse d'un complexe non nul}

\textbf{Énoncé.} Pour tout $z\neq 0$,
\[
\frac{1}{z}=\frac{\overline z}{|z|^2}.
\]

\textbf{Type de raisonnement.} Vérification directe de la propriété caractéristique de l'inverse.

\begin{prop}
Soit $z\neq 0$. Alors
\[
|z|^2=z\overline z\neq 0,
\]
de sorte que le nombre $\dfrac{\overline z}{|z|^2}$ est bien défini.

Calculons :
\[
z\cdot \frac{\overline z}{|z|^2}
=\frac{z\overline z}{|z|^2}
=\frac{|z|^2}{|z|^2}
=1.
\]
Ainsi, le nombre $\dfrac{\overline z}{|z|^2}$ multiplié par $z$ donne $1$ ; c'est donc l'inverse de $z$.

On en déduit
\[
\frac{1}{z}=\frac{\overline z}{|z|^2}.
\]
\end{prop}

\section*{7. Si \texorpdfstring{$z=r(\cos\theta+\ii\sin\theta)$}{z=r(cosθ+i sinθ)}, alors \texorpdfstring{$|z|=r$}{|z|=r}}

\textbf{Énoncé.} Si
\[
z=r(\cos\theta+\ii\sin\theta)
\qquad \text{avec } r\geq 0,
\]
alors
\[
|z|=r.
\]

\textbf{Type de raisonnement.} Calcul direct utilisant l'identité trigonométrique fondamentale.

\begin{prop}
On écrit
\[
z=(r\cos\theta)+\ii(r\sin\theta).
\]
Ainsi,
\[
\Reel(z)=r\cos\theta,
\qquad
\Imag(z)=r\sin\theta.
\]
Par définition du module,
\[
|z|=\sqrt{(r\cos\theta)^2+(r\sin\theta)^2}.
\]
On factorise par $r^2$ :
\[
|z|=\sqrt{r^2(\cos^2\theta+\sin^2\theta)}.
\]
Or
\[
\cos^2\theta+\sin^2\theta=1.
\]
Donc
\[
|z|=\sqrt{r^2}=r,
\]
puisque $r\geq 0$.
\end{prop}

\section*{8. Affixe du milieu d'un segment}

\textbf{Énoncé.} Soient $A$ et $B$ d'affixes respectives $z_A$ et $z_B$. Le milieu $M$ de $[AB]$ a pour affixe
\[
z_M=\frac{z_A+z_B}{2}.
\]

\textbf{Type de raisonnement.} Traduction complexe d'une relation vectorielle caractéristique.

\begin{prop}
Par définition du milieu, le point $M$ vérifie
\[
\overrightarrow{AM}=\frac12\,\overrightarrow{AB}.
\]
En écriture complexe, cela se traduit par
\[
z_M-z_A=\frac12(z_B-z_A).
\]
En développant le membre de droite,
\[
z_M-z_A=\frac12 z_B-\frac12 z_A.
\]
Il suffit alors d'ajouter $z_A$ aux deux membres :
\[
z_M=\frac12 z_B+\frac12 z_A=\frac{z_A+z_B}{2}.
\]

On retrouve bien l'affixe du milieu.
\end{prop}

\section*{9. Condition pour appartenir à un cercle}

\textbf{Énoncé.} Soit $A$ d'affixe $a$. Un point $M$ d'affixe $z$ appartient au cercle de centre $A$ et de rayon $R$ si et seulement si
\[
|z-a|=R.
\]

\textbf{Type de raisonnement.} Caractérisation géométrique traduite en écriture complexe.

\begin{prop}
L'affixe du vecteur $\overrightarrow{AM}$ est
\[
z-a.
\]
Le module de ce complexe est donc la longueur du vecteur $\overrightarrow{AM}$, c'est-à-dire
\[
|z-a|=AM.
\]

Par définition, le point $M$ appartient au cercle de centre $A$ et de rayon $R$ si et seulement si
\[
AM=R.
\]
En remplaçant $AM$ par $|z-a|$, on obtient
\[
M\in \mathcal{C}(A,R)\iff |z-a|=R.
\]
\end{prop}

\section*{10. Alignement de trois points}

\textbf{Énoncé.} Soient $A,B,C$ d'affixes respectives $z_A,z_B,z_C$, avec $z_B\neq z_A$. Alors
\[
A,B,C \text{ alignés } \iff \frac{z_C-z_A}{z_B-z_A}\in\R.
\]

\textbf{Type de raisonnement.} Double implication fondée sur la caractérisation de la colinéarité.

\begin{prop}
Les points $A,B,C$ sont alignés si et seulement si les vecteurs $\overrightarrow{AB}$ et $\overrightarrow{AC}$ sont colinéaires.

Or l'affixe de $\overrightarrow{AB}$ est
\[
z_B-z_A,
\]
et celle de $\overrightarrow{AC}$ est
\[
z_C-z_A.
\]

Les vecteurs $\overrightarrow{AB}$ et $\overrightarrow{AC}$ sont colinéaires si et seulement s'il existe un réel $k$ tel que
\[
z_C-z_A=k(z_B-z_A).
\]
Comme $z_B-z_A\neq 0$, cette relation équivaut à
\[
\frac{z_C-z_A}{z_B-z_A}=k.
\]
Comme $k$ est réel, on obtient
\[
\frac{z_C-z_A}{z_B-z_A}\in\R.
\]

Réciproquement, si
\[
\frac{z_C-z_A}{z_B-z_A}\in\R,
\]
alors il existe $k\in\R$ tel que
\[
\frac{z_C-z_A}{z_B-z_A}=k,
\]
c'est-à-dire
\[
z_C-z_A=k(z_B-z_A).
\]
Les vecteurs $\overrightarrow{AC}$ et $\overrightarrow{AB}$ sont donc colinéaires, ce qui prouve que les points $A,B,C$ sont alignés.
\end{prop}

\section*{11. Orthogonalité de deux vecteurs par quotient imaginaire pur}

\textbf{Énoncé.} Soient $A,B,C$ trois points tels que $z_B\neq z_A$ et $z_C\neq z_A$. Alors
\[
\overrightarrow{AB}\perp \overrightarrow{AC}
\iff
\frac{z_C-z_A}{z_B-z_A}\in \ii\R.
\]

\textbf{Type de raisonnement.} Double implication fondée sur la caractérisation complexe d'un quart de tour.

\begin{prop}
Supposons d'abord que
\[
\overrightarrow{AB}\perp \overrightarrow{AC}.
\]
Deux vecteurs non nuls sont orthogonaux si et seulement si l'un s'obtient à partir de l'autre par multiplication par un imaginaire pur. Il existe donc $\lambda\in\R$ tel que
\[
\overrightarrow{AC}=\ii\lambda\,\overrightarrow{AB}.
\]
En écriture complexe, cela donne
\[
z_C-z_A=\ii\lambda(z_B-z_A).
\]
Comme $z_B-z_A\neq 0$, on peut diviser par $z_B-z_A$ :
\[
\frac{z_C-z_A}{z_B-z_A}=\ii\lambda.
\]
Or $\ii\lambda\in \ii\R$, donc
\[
\frac{z_C-z_A}{z_B-z_A}\in \ii\R.
\]

Réciproquement, supposons que
\[
\frac{z_C-z_A}{z_B-z_A}\in \ii\R.
\]
Alors il existe $\lambda\in\R$ tel que
\[
\frac{z_C-z_A}{z_B-z_A}=\ii\lambda.
\]
En multipliant par $z_B-z_A$, on obtient
\[
z_C-z_A=\ii\lambda(z_B-z_A).
\]
Le vecteur $\overrightarrow{AC}$ s'obtient donc à partir de $\overrightarrow{AB}$ par multiplication par un imaginaire pur ; il lui est donc orthogonal. Ainsi
\[
\overrightarrow{AB}\perp \overrightarrow{AC}.
\]
\end{prop}

\section*{12. Caractérisation d'une médiatrice par les modules}

\textbf{Énoncé.} Soient $A$ et $B$ d'affixes $a$ et $b$. Un point $M$ d'affixe $z$ appartient à la médiatrice de $[AB]$ si et seulement si
\[
|z-a|=|z-b|.
\]

\textbf{Type de raisonnement.} Double implication par traduction d'une caractérisation géométrique classique.

\begin{prop}
Par définition, un point appartient à la médiatrice de $[AB]$ si et seulement s'il est à égale distance de $A$ et de $B$.

Or, si $M$ est d'affixe $z$, on a
\[
AM=|z-a|
\qquad \text{et} \qquad
BM=|z-b|.
\]
Par conséquent,
\[
M \text{ appartient à la médiatrice de } [AB]
\iff
AM=BM
\iff
|z-a|=|z-b|.
\]
\end{prop}

\section*{13. Cercle de diamètre \texorpdfstring{$[AB]$}{[AB]} : angle droit}

\textbf{Énoncé.} Soient $A$ et $B$ deux points distincts d'affixes $a$ et $b$. Pour un point $M$ d'affixe $z$, on a
\[
\widehat{AMB}=90^\circ
\iff
\frac{z-a}{z-b}\in \ii\R.
\]

\textbf{Type de raisonnement.} Réduction à une condition d'orthogonalité, puis utilisation d'une caractérisation précédente.

\begin{prop}
L'affixe du vecteur $\overrightarrow{MA}$ est
\[
a-z=-(z-a),
\]
et celle du vecteur $\overrightarrow{MB}$ est
\[
b-z=-(z-b).
\]
On en déduit
\[
\frac{a-z}{b-z}=\frac{-(z-a)}{-(z-b)}=\frac{z-a}{z-b}.
\]

Or
\[
\widehat{AMB}=90^\circ
\iff
\overrightarrow{MA}\perp \overrightarrow{MB}.
\]
D'après la caractérisation de l'orthogonalité obtenue précédemment, cela équivaut à
\[
\frac{a-z}{b-z}\in \ii\R.
\]
Comme
\[
\frac{a-z}{b-z}=\frac{z-a}{z-b},
\]
on obtient finalement
\[
\widehat{AMB}=90^\circ
\iff
\frac{z-a}{z-b}\in \ii\R.
\]
\end{prop}

\section*{14. Ensemble des points \texorpdfstring{$M$}{M} tels que \texorpdfstring{$\left|\dfrac{z-a}{z-b}\right|=1$}{|(z-a)/(z-b)|=1}}

\textbf{Énoncé.} Soient $A$ et $B$ deux points distincts d'affixes $a$ et $b$. L'ensemble des points $M$ d'affixe $z$, avec $z\neq b$, tels que
\[
\left|\frac{z-a}{z-b}\right|=1
\]
est la médiatrice de $[AB]$.

\textbf{Type de raisonnement.} Transformation algébrique de la condition, puis interprétation géométrique.

\begin{prop}
On utilise la formule du module d'un quotient :
\[
\left|\frac{z-a}{z-b}\right|=\frac{|z-a|}{|z-b|}.
\]
Ainsi,
\[
\left|\frac{z-a}{z-b}\right|=1
\iff
\frac{|z-a|}{|z-b|}=1.
\]
Comme $|z-b|>0$ puisque $z\neq b$, cette relation équivaut à
\[
|z-a|=|z-b|.
\]

Or cette dernière égalité signifie que le point $M$ est à égale distance de $A$ et de $B$. L'ensemble des points vérifiant cette propriété est précisément la médiatrice de $[AB]$.

L'ensemble cherché est donc la médiatrice de $[AB]$.
\end{prop}

\section*{15. Nature d'une transformation \texorpdfstring{$z'=az+b$}{z'=az+b}}

\textbf{Énoncé.} Soient $a,b\in\C$ avec $a\neq 0$. L'application
\[
z'=az+b
\]
transforme toute droite en droite et tout cercle en cercle.

\textbf{Type de raisonnement.} Décomposition en transformations élémentaires et composition de propriétés géométriques.

\begin{prop}
On peut réécrire l'application sous la forme
\[
z'=a\left(z+\frac{b}{a}\right).
\]
Elle apparaît ainsi comme la composée de deux transformations :
\begin{itemize}[label=--]
    \item la translation $z\mapsto z+\dfrac{b}{a}$ ;
    \item puis la multiplication par $a$.
\end{itemize}

La translation conserve les alignements et les distances entre points ; elle transforme donc une droite en une droite et un cercle en un cercle.

Étudions maintenant la multiplication par $a$. Écrivons
\[
a=\rho(\cos\theta+\ii\sin\theta)
\qquad \text{avec } \rho=|a|>0.
\]
Multiplier un complexe par $a$, c'est :
\begin{itemize}[label=--]
    \item multiplier son module par $\rho$ ;
    \item ajouter $\theta$ à son argument.
\end{itemize}
Autrement dit, cette transformation est la composée d'une homothétie de centre $O$ et de rapport $\rho$, puis d'une rotation de centre $O$ et d'angle $\theta$.

Or une homothétie transforme une droite en une droite et un cercle en un cercle. Il en est de même d'une rotation.

La composée de transformations qui envoient les droites sur des droites et les cercles sur des cercles possède encore cette propriété. On en déduit que l'application
\[
z\mapsto az+b
\]
transforme toute droite en droite et tout cercle en cercle.
\end{prop}

\vspace{1em}
\hrule
\vspace{1em}

\section*{Pour apprendre efficacement}

On peut découper ce programme de colle en trois blocs.

\begin{itemize}
    \item \textbf{Bloc 1 --- Calcul algébrique :} démonstrations 1 à 7.
    \item \textbf{Bloc 2 --- Géométrie complexe de base :} démonstrations 8 à 10.
    \item \textbf{Bloc 3 --- Géométrie complexe plus ambitieuse :} démonstrations 11 à 15.
\end{itemize}

\end{document}